\begin{PSbox}[unbreakable]{obj}{Learning Objectives}

After completing this module, students will be able to:

\begin{enumerate}
\def\labelenumi{\arabic{enumi}.}
\tightlist
\item
  Compute expected values for discrete and continuous random variables
\item
  Calculate variance and standard deviation and interpret their engineering meaning
\item
  Apply properties of expectation and variance in engineering analysis
\item
  Compute and interpret higher-order moments
\item
  Use the moment generating function (MGF)
\item
  Relate statistical moments to physical engineering quantities
\end{enumerate}

\end{PSbox}

\PSsection{4.1}{Introduction: Why Summary Statistics?}

\PSsubsection{The Engineering Need}

A probability distribution contains ALL information about a random variable. But for engineering decisions, we often need concise summaries:

\begin{itemize}
\tightlist
\item
  \textbf{Signal processing:} What is the average power? \ensuremath{\to} Mean of $X^{2}$
\item
  \textbf{Quality control:} How much $do$ components vary? \ensuremath{\to} Variance
\item
  \textbf{Communication:} What is the noise power? $\to E[N^{2}]$
\item
  \textbf{Reliability:} What is the average lifetime? \ensuremath{\to} Mean
\end{itemize}

Moments (mean, variance, etc.) are the essential \textbf{engineering-relevant summaries} of a distribution.

\PSsection{4.2}{Expected Value (Mean)}

\begin{PSbox}[unbreakable]{def}{Definition — Discrete RV}

\PSdisplay{E[X] = \mu_X = \sum_{x} x \cdot p_X(x)}

Sum of (each value \ensuremath{\times} its probability).

\end{PSbox}

\begin{PSbox}[unbreakable]{def}{Definition — Continuous RV}

\PSdisplay{E[X] = \mu_X = \int_{-\infty}^{\infty} x \cdot f_X(x) \, dx}

\end{PSbox}

\PSsubsection{Engineering Interpretation}

The expected value is the \textbf{long-run average}: if you repeat the experiment many times and average the results, you get approximately $E[X]$.

{\def\LTcaptype{none} % do not increment counter
\begin{longtable}[c]{@{}cc@{}}
\toprule\noalign{}
\rowcolor{PSnavy}\PSth{Engineering Quantity} & \PSth{Expected Value Meaning} \\
\midrule\noalign{}
\endhead
\bottomrule\noalign{}
\endlastfoot
Noise voltage $E[N]$ & DC bias/offset in noise \\
Signal amplitude $E[A]$ & Average signal level \\
Component lifetime $E[T]$ & Mean time to failure (MTTF) \\
Packet delay $E[D]$ & Average delay \\
Bit errors $E[X]$ & Average error count \\
\end{longtable}
}

\begin{PSbox}[unbreakable]{ex}{Example: Average Noise Power}

For zero-mean Gaussian noise $N \sim N(0,\, \sigma^{2})$:

\begin{itemize}
\tightlist
\item
  $E[N] = 0$ (no DC component)
\item
  Average power $= E[N^{2}] = \sigma^{2}$ (this is the variance!)
\end{itemize}

\end{PSbox}

\begin{PSbox}[unbreakable]{ex}{Example: Expected Number of Retransmissions}

ARQ protocol: each transmission succeeds with probability $p = 0.8$.\PSbr{}Number of attempts $X \sim \mathrm{Geometric}(0.8)$.\PSbr{}$E[X] = \frac{1}{p} = \frac{1}{0.8} = 1.25$ attempts on average.

\end{PSbox}

\PSsection{4.3}{Expected Value of a Function: $E[g(X)]$}

\PSsubsection{The LOTUS Theorem (Law of the Unconscious Statistician)}

To find $E[g(X)]$, you $do$ NOT need the distribution of $Y = g(X)$:

\textbf{Discrete:}\PSdisplay{E[g(X)] = \sum_{x} g(x) \cdot p_X(x)}

\textbf{Continuous:}\PSdisplay{E[g(X)] = \int_{-\infty}^{\infty} g(x) \cdot f_X(x) \, dx}

\PSsubsection{Engineering Applications}

{\def\LTcaptype{none} % do not increment counter
\begin{longtable}[c]{@{}cc@{}}
\toprule\noalign{}
\rowcolor{PSnavy}\PSth{Function $g(X)$} & \PSth{Engineering Meaning} \\
\midrule\noalign{}
\endhead
\bottomrule\noalign{}
\endlastfoot
$g(X) = X^{2}$ & Mean square value (power) \\
$g(X) = (X - \mu)^{2}$ & Variance (AC power) \\
$g(X)$ = & $X$ \\
$g(X) = e^{j\omega X}$ & Characteristic function \\
\end{longtable}
}

\begin{PSbox}[unbreakable]{ex}{Example: Mean Power of a Signal}

Signal $X$ has PDF $f_{X}(x)$. Power $P = X^{2}$.

$E[P] = E[X^{2}] = \int x^{2} f_{X}(x)\,dx$

For $X \sim N(0,\, \sigma^{2})$: $E[X^{2}] = \sigma^{2}$ (the noise power equals the variance).

For $X \sim N(\mu,\, \sigma^{2})$: $E[X^{2}] = \mu^{2} + \sigma^{2}$ (total power = DC power + AC power).

\end{PSbox}

\PSsection{4.4}{Variance}

\begin{PSbox}[unbreakable]{def}{Definition}

\PSdisplay{\text{Var}(X) = E[(X - \mu_X)^2] = E[X^2] - (E[X])^2}

The second formula (computational formula) is often easier to use.

\end{PSbox}

\PSsubsection{Derivation of Computational Formula}

$\operatorname{Var}(X) = E[(X - \mu)^{2}]$\PSbr{}$= E[X^{2} - 2\mu X + \mu^{2}]$\PSbr{}$= E[X^{2}] - 2\mu E[X] + \mu^{2}$\PSbr{}$= E[X^{2}] - 2\mu^{2} + \mu^{2}$\PSbr{}$= E[X^{2}] - \mu^{2}$

\PSsubsection{Engineering Interpretation}

The variance measures \textbf{spread/dispersion/uncertainty}:

{\def\LTcaptype{none} % do not increment counter
\begin{longtable}[c]{@{}cc@{}}
\toprule\noalign{}
\rowcolor{PSnavy}\PSth{Context} & \PSth{Variance Meaning} \\
\midrule\noalign{}
\endhead
\bottomrule\noalign{}
\endlastfoot
Noise signal & Noise power (for zero-mean noise) \\
Manufacturing & How much components deviate from nominal \\
Measurements & Measurement precision (smaller = more precise) \\
Signal & AC power component \\
Lifetime & How predictable is the lifetime? \\
\end{longtable}
}

\PSsubsection{Key Insight: Variance = Power in Engineering}

For a zero-mean signal $X$:

\begin{itemize}
\tightlist
\item
  $E[X^{2}] = \operatorname{Var}(X) = \sigma^{2}$ = \textbf{total power}
\item
  For $X \sim N(0,\, \sigma^{2})$: noise power $= \sigma^{2} = \operatorname{Var}(X)$
\end{itemize}

For a signal with DC offset $(\mu \ne 0)$:

\begin{itemize}
\tightlist
\item
  $E[X^{2}] = \mu^{2} + \sigma^{2}$ = \textbf{total power = DC power + AC power}
\item
  $\operatorname{Var}(X) = \sigma^{2}$ = \textbf{AC power only}
\end{itemize}

\PSsection{4.5}{Standard Deviation}

\begin{PSbox}[unbreakable]{def}{Definition}

\PSdisplay{\sigma_X = \sqrt{\text{Var}(X)}}

\end{PSbox}

\PSsubsection{Engineering Significance}

The standard deviation has the \textbf{same units} as $X$, making it directly interpretable:

{\def\LTcaptype{none} % do not increment counter
\begin{longtable}[c]{@{}ccc@{}}
\toprule\noalign{}
\rowcolor{PSnavy}\PSth{Random Variable} & \PSth{Standard Deviation Unit} & \PSth{Interpretation} \\
\midrule\noalign{}
\endhead
\bottomrule\noalign{}
\endlastfoot
Voltage noise & Volts & RMS noise voltage \\
Timing jitter & Seconds & RMS jitter \\
Distance error & Meters & RMS position error \\
Temperature & \ensuremath{^{\circ}}C & RMS temperature fluctuation \\
\end{longtable}
}

The standard deviation is the \textbf{RMS (root-mean-square)} value for zero-mean signals.

\PSsection{4.6}{Properties of Expectation}

\PSsubsection{Linearity (FUNDAMENTAL)}

\PSdisplay{E[aX + b] = aE[X] + b}

\PSdisplay{E[X + Y] = E[X] + E[Y] \quad \text{(ALWAYS, even if dependent!)}}

\PSdisplay{E\left[\sum_{i=1}^n a_i X_i\right] = \sum_{i=1}^n a_i E[X_i]}

\PSsubsection{Monotonicity}

If $X \ge 0$ always, then $E[X] \ge 0$.

If $X \ge Y$ always, then $E[X] \ge E[Y]$.

\PSsubsection{For Independent RVs}

If $X$ and $Y$ are independent:\PSdisplay{E[XY] = E[X] \cdot E[Y]}

\begin{PSquote}

\PSwarn{} This does NOT hold for dependent variables!

\end{PSquote}

\begin{PSbox}[unbreakable]{ex}{Engineering Example: Total Noise}

A received signal: $R = S + N_{1} + N_{2}$ (signal plus two noise sources)

$E[R] = E[S] + E[N_{1}] + E[N_{2}] = E[S] + 0 + 0 = E[S]$

(Zero-mean noise doesn’t bias the expected received signal.)

\end{PSbox}

\PSsection{4.7}{Properties of Variance}

\PSsubsection{Scaling and Shift}

\PSdisplay{\text{Var}(aX + b) = a^2 \text{Var}(X)}

\begin{itemize}
\tightlist
\item
  Adding a constant $b$ does NOT change variance (shifting doesn’t change spread)
\item
  Scaling by a multiplies variance by $a^{2}$
\end{itemize}

\PSsubsection{Sum of Independent RVs}

If $X_{1},\, X_{2},\, \ldots,\, X_{n}$ are \textbf{independent}:

\PSdisplay{\text{Var}(X_1 + X_2 + \cdots + X_n) = \text{Var}(X_1) + \text{Var}(X_2) + \cdots + \text{Var}(X_n)}

\begin{PSquote}

\PSwarn{} Variances add for INDEPENDENT variables only!

\end{PSquote}

\begin{PSbox}[unbreakable]{ex}{Engineering Example: Noise Power Addition}

Two independent noise sources $(\sigma_{1}^{2} = 0.01\,\mathrm{V}^{2},\, \sigma_{2}^{2} = 0.04\,\mathrm{V}^{2})$:

Total noise power $= \sigma_{1}^{2} + \sigma_{2}^{2} = 0.05\,\mathrm{V}^{2}$\PSbr{}Total RMS noise $= \sqrt{0.05} = 0.224\,\mathrm{V}$

Note: RMS values $do$ NOT add directly! $(0.1 + 0.2 \ne 0.224)$

\end{PSbox}

\PSsubsection{Variance of Sample Mean}

For $n$ i.i.d. observations $X_{1},\, \ldots,\, X_{n}$ with variance $\sigma^{2}$:

\PSdisplay{\text{Var}(\bar{X}) = \text{Var}\left(\frac{1}{n}\sum X_i\right) = \frac{\sigma^2}{n}}

More measurements \ensuremath{\to} lower uncertainty in the average.

\PSsection{4.8}{Higher-Order Moments}

\PSsubsection{$n-\mathrm{th}$ Moment}

\PSdisplay{\mu_n' = E[X^n] = \int_{-\infty}^{\infty} x^n f_X(x) \, dx}

\PSsubsection{$n-\mathrm{th}$ Central Moment}

\PSdisplay{\mu_n = E[(X - \mu)^n]}

\begin{itemize}
\tightlist
\item
  $\mu_{1} = 0$ (always)
\item
  $\mu_{2} = \operatorname{Var}(X)$
\item
  $\mu_{3}$ relates to \textbf{skewness}
\item
  $\mu_{4}$ relates to \textbf{kurtosis}
\end{itemize}

\PSsubsection{Skewness ($3rd$ Central Moment, Normalized)}

\PSdisplay{\gamma_1 = \frac{E[(X-\mu)^3]}{\sigma^3}}

{\def\LTcaptype{none} % do not increment counter
\begin{longtable}[c]{@{}
  >{\centering\arraybackslash}p{(0.965\linewidth - 4\tabcolsep) * \real{0.3846}}
  >{\centering\arraybackslash}p{(0.965\linewidth - 4\tabcolsep) * \real{0.2692}}
  >{\centering\arraybackslash}p{(0.965\linewidth - 4\tabcolsep) * \real{0.3462}}@{}}
\toprule\noalign{}
\rowcolor{PSnavy}\PSth{Skewness} & \PSth{Shape} & \PSth{Example} \\
\midrule\noalign{}
\endhead
\bottomrule\noalign{}
\endlastfoot
$\gamma_{1} = 0$ & Symmetric & Gaussian \\
$\gamma_{1} > 0$ & Right-skewed (long right tail) & Exponential \\
$\gamma_{1} < 0$ & Left-skewed (long left tail) & Negated exponential \\
\end{longtable}
}

\textbf{Engineering significance:} Skewness tells us if extreme values are more likely in one direction. Important for reliability (failure time distributions are right-skewed).

\PSsubsection{Kurtosis ($4\mathrm{th}$ Central Moment, Normalized)}

\PSdisplay{\gamma_2 = \frac{E[(X-\mu)^4]}{\sigma^4} - 3}

(Subtracting 3 makes Gaussian have kurtosis 0: ``excess kurtosis.'')

{\def\LTcaptype{none} % do not increment counter
\begin{longtable}[c]{@{}ccc@{}}
\toprule\noalign{}
\rowcolor{PSnavy}\PSth{Kurtosis} & \PSth{Interpretation} & \PSth{Meaning} \\
\midrule\noalign{}
\endhead
\bottomrule\noalign{}
\endlastfoot
$\gamma_{2} = 0$ & Mesokurtic & Gaussian-like tails \\
$\gamma_{2} > 0$ & Leptokurtic & Heavier tails than Gaussian \\
$\gamma_{2} < 0$ & Platykurtic & Lighter tails than Gaussian \\
\end{longtable}
}

\textbf{Engineering significance:} High kurtosis means more outliers — important for impulsive noise modeling.

\PSsection{4.9}{Moment Generating Function (MGF)}

\begin{PSbox}[unbreakable]{def}{Definition}

\PSdisplay{M_X(t) = E[e^{tX}]}

\textbf{Discrete:} $M_{X}(t) = \sum e^{tx} p_{X}(x)$

\textbf{Continuous:} $M_{X}(t) = \int e^{tx} f_{X}(x)\,dx$

\end{PSbox}

\PSsubsection{Why ``Moment Generating''?}

The $n-\mathrm{th}$ moment is obtained by differentiating $n$ times at $t = 0$:

\PSdisplay{E[X^n] = \frac{d^n M_X(t)}{dt^n}\bigg|_{t=0}}

\begin{itemize}
\tightlist
\item
  $M_{X}'(0) = E[X]$
\item
  $M_{X}''(0) = E[X^{2}]$
\item
  $\operatorname{Var}(X) = M_{X}''(0) - [M_{X}'(0)]^{2}$
\end{itemize}

\PSsubsection{Properties}

\begin{enumerate}
\def\labelenumi{\arabic{enumi}.}
\item
  \textbf{Uniqueness:} If $M_{X}(t) = M_{Y}(t)$ for all $t$ in a neighborhood of 0, then $X$ and $Y$ have the same distribution.
\item
  \textbf{Sum of Independent RVs:} If $X$ and $Y$ are independent:\PSbr{}$M_{X+Y}(t) = M_{X}(t) \cdot M_{Y}(t)$
\item
  \textbf{Linear transformation:} $M_{\text{aX}+b}(t) = e^{bt} \cdot M_{X}(at)$
\end{enumerate}

\PSsubsection{MGFs of Common Distributions}

{\def\LTcaptype{none} % do not increment counter
\begin{longtable}[c]{@{}
  >{\centering\arraybackslash}p{(0.965\linewidth - 2\tabcolsep) * \real{0.5200}}
  >{\centering\arraybackslash}p{(0.965\linewidth - 2\tabcolsep) * \real{0.4800}}@{}}
\toprule\noalign{}
\rowcolor{PSnavy}\PSth{Distribution} & \PSth{MGF $M_{X}(t)$} \\
\midrule\noalign{}
\endhead
\bottomrule\noalign{}
\endlastfoot
$\mathrm{Bernoulli}(p)$ & $1 - p + pe^{t}$ \\
$\mathrm{Binomial}(n,\,p)$ & $(1 - p + pe^{t})^{n}$ \\
$\mathrm{Poisson}(\lambda)$ & $\exp (\lambda (e^{t} - 1))$ \\
$\mathrm{Exponential}(\lambda)$ & $\frac{\lambda}{\lambda - t},\, t < \lambda$ \\
$\mathrm{Gaussian}(\mu,\,\sigma^{2})$ & $\exp \left(\mu t + \frac{\sigma^{2}t^{2}}{2}\right)$ \\
\end{longtable}
}

\PSsubsection{Engineering Application: Sum of Independent Signals}

If signal components $X_{1} \sim N(\mu_{1},\, \sigma_{1}^{2})$ and $X_{2} \sim N(\mu_{2},\, \sigma_{2}^{2})$ are independent:

$M_{X_{1}+X_{2}}(t) = M_{X_{1}}(t) \cdot M_{X_{2}}(t)$\PSbr{}$\displaystyle = \exp \left(\mu_{1}t + \frac{\sigma_{1}^{2}t^{2}}{2}\right) \cdot \exp \left(\mu_{2}t + \frac{\sigma_{2}^{2}t^{2}}{2}\right)$\PSbr{}$\displaystyle = \exp \left((\mu_{1}+\mu_{2})t + \frac{(\sigma_{1}^{2}+\sigma_{2}^{2})t^{2}}{2}\right)$

This is the MGF of $N(\mu_{1}+\mu_{2},\, \sigma_{1}^{2}+\sigma_{2}^{2})$. Therefore $X_{1}+X_{2} \sim N(\mu_{1}+\mu_{2},\, \sigma_{1}^{2}+\sigma_{2}^{2})$.

\PSsection{4.10}{Engineering Applications of Moments}

\PSsubsection{Application 1: Noise Power Characterization}

A noise signal $N(t)$ is sampled. The samples have mean $\mu_{N}$ and variance $\sigma_{N}^{2}$.

\begin{itemize}
\tightlist
\item
  \textbf{DC component (bias):} $\mu_{N}$ — should be zero for ideal noise
\item
  \textbf{AC noise power:} $\sigma_{N}^{2}$ — determines SNR
\item
  \textbf{Total power:} $E[N^{2}] = \mu_{N}^{2} + \sigma_{N}^{2}$
\end{itemize}

\textbf{SNR (Signal-to-Noise Ratio):}\PSdisplay{\text{SNR} = \frac{E[S^2]}{\sigma_N^2} = \frac{\text{Signal Power}}{\text{Noise Power}}}

\PSsubsection{Application 2: Measurement Error}

A sensor measures true value $\theta$ with error $E$:

\begin{itemize}
\tightlist
\item
  Measurement: $X = \theta + E$
\item
  $E[X] = \theta + E[E]$ — bias in measurement
\item
  $\operatorname{Var}(X) = \operatorname{Var}(E)$ — measurement precision
\end{itemize}

\textbf{Accuracy:} $\lvert E[E]\rvert$ (how close the average is to truth)\PSbr{}\textbf{Precision:} $\sigma_{E}$ (how repeatable measurements are)

\PSsubsection{Application 3: Component Lifetime Statistics}

Components have lifetime $T \sim \mathrm{Exponential}(\lambda)$:

\begin{itemize}
\tightlist
\item
  Mean Time To Failure: $\mathrm{MTTF} = E[T] = \frac{1}{\lambda}$
\item
  Standard Deviation: $\sigma_{T} = \frac{1}{\lambda}$ = MTTF
\item
  For exponential: coefficient of variation CV $= \frac{\sigma}{\mu} = 1$ (always!)
\end{itemize}

\PSsubsection{Application 4: Communication Signal Amplitude}

Received signal $R = A + N$ where $A$ = signal amplitude, $N \sim N(0,\, \sigma^{2})$:

\begin{itemize}
\tightlist
\item
  $E[R] = A$ (signal amplitude preserved on average)
\item
  $\operatorname{Var}(R) = \sigma^{2}$ (noise power determines variance)
\item
  $\displaystyle \mathrm{SNR} = \frac{A^{2}}{\sigma^{2}}$
\end{itemize}

\PSsection{4.11}{MATLAB Examples}

\begin{PSbox}[unbreakable]{ex}{Example 1: Computing Moments from a Distribution}

\tcbset{PScodemode/.style={unbreakable}}

\begin{Shaded}
\begin{Highlighting}[]
\CommentTok{\%\% Moments of Common Distributions}

\CommentTok{\% Exponential: failure time, lambda = 0.01 (per hour)}
\VariableTok{lambda} \OperatorTok{=} \FloatTok{0.01}\OperatorTok{;}
\VariableTok{beta} \OperatorTok{=} \FloatTok{1}\OperatorTok{/}\VariableTok{lambda}\OperatorTok{;}  \CommentTok{\% mean = 100 hours}

\VariableTok{fprintf}\NormalTok{(}\SpecialStringTok{\textquotesingle{}=== Exponential Distribution (λ = \%.3f) ===\textbackslash{}n\textquotesingle{}}\OperatorTok{,} \VariableTok{lambda}\NormalTok{)}\OperatorTok{;}
\VariableTok{fprintf}\NormalTok{(}\SpecialStringTok{\textquotesingle{}Mean (MTTF) = \%.1f hours\textbackslash{}n\textquotesingle{}}\OperatorTok{,} \VariableTok{beta}\NormalTok{)}\OperatorTok{;}
\VariableTok{fprintf}\NormalTok{(}\SpecialStringTok{\textquotesingle{}Variance = \%.1f hours²\textbackslash{}n\textquotesingle{}}\OperatorTok{,} \VariableTok{beta}\OperatorTok{\^{}}\FloatTok{2}\NormalTok{)}\OperatorTok{;}
\VariableTok{fprintf}\NormalTok{(}\SpecialStringTok{\textquotesingle{}Std Dev = \%.1f hours\textbackslash{}n\textquotesingle{}}\OperatorTok{,} \VariableTok{beta}\NormalTok{)}\OperatorTok{;}
\VariableTok{fprintf}\NormalTok{(}\SpecialStringTok{\textquotesingle{}Coefficient of Variation = \%.2f\textbackslash{}n\textquotesingle{}}\OperatorTok{,} \VariableTok{beta}\OperatorTok{/}\VariableTok{beta}\NormalTok{)}\OperatorTok{;}

\CommentTok{\% Gaussian: noise voltage, sigma = 0.1V}
\VariableTok{mu} \OperatorTok{=} \FloatTok{0}\OperatorTok{;} \VariableTok{sigma} \OperatorTok{=} \FloatTok{0.1}\OperatorTok{;}
\VariableTok{fprintf}\NormalTok{(}\SpecialStringTok{\textquotesingle{}\textbackslash{}n=== Gaussian Distribution (μ=\%.1f, σ=\%.2f V) ===\textbackslash{}n\textquotesingle{}}\OperatorTok{,} \VariableTok{mu}\OperatorTok{,} \VariableTok{sigma}\NormalTok{)}\OperatorTok{;}
\VariableTok{fprintf}\NormalTok{(}\SpecialStringTok{\textquotesingle{}Mean = \%.2f V\textbackslash{}n\textquotesingle{}}\OperatorTok{,} \VariableTok{mu}\NormalTok{)}\OperatorTok{;}
\VariableTok{fprintf}\NormalTok{(}\SpecialStringTok{\textquotesingle{}Variance (noise power) = \%.4f V²\textbackslash{}n\textquotesingle{}}\OperatorTok{,} \VariableTok{sigma}\OperatorTok{\^{}}\FloatTok{2}\NormalTok{)}\OperatorTok{;}
\VariableTok{fprintf}\NormalTok{(}\SpecialStringTok{\textquotesingle{}RMS noise = \%.3f V\textbackslash{}n\textquotesingle{}}\OperatorTok{,} \VariableTok{sigma}\NormalTok{)}\OperatorTok{;}
\VariableTok{fprintf}\NormalTok{(}\SpecialStringTok{\textquotesingle{}E[X²] (total power) = \%.4f V²\textbackslash{}n\textquotesingle{}}\OperatorTok{,} \VariableTok{mu}\OperatorTok{\^{}}\FloatTok{2} \OperatorTok{+} \VariableTok{sigma}\OperatorTok{\^{}}\FloatTok{2}\NormalTok{)}\OperatorTok{;}
\end{Highlighting}
\end{Shaded}

\end{PSbox}

\begin{PSbox}[unbreakable]{ex}{Example 2: Verifying Variance Properties}

\tcbset{PScodemode/.style={unbreakable}}

\begin{Shaded}
\begin{Highlighting}[]
\CommentTok{\%\% Variance of Sum of Independent RVs}
\VariableTok{N} \OperatorTok{=} \FloatTok{100000}\OperatorTok{;}

\CommentTok{\% Two independent noise sources}
\VariableTok{sigma1} \OperatorTok{=} \FloatTok{0.1}\OperatorTok{;}   \CommentTok{\% RMS of noise 1}
\VariableTok{sigma2} \OperatorTok{=} \FloatTok{0.2}\OperatorTok{;}   \CommentTok{\% RMS of noise 2}

\VariableTok{X1} \OperatorTok{=} \VariableTok{sigma1} \OperatorTok{*} \VariableTok{randn}\NormalTok{(}\FloatTok{1}\OperatorTok{,} \VariableTok{N}\NormalTok{)}\OperatorTok{;}
\VariableTok{X2} \OperatorTok{=} \VariableTok{sigma2} \OperatorTok{*} \VariableTok{randn}\NormalTok{(}\FloatTok{1}\OperatorTok{,} \VariableTok{N}\NormalTok{)}\OperatorTok{;}
\VariableTok{Y} \OperatorTok{=} \VariableTok{X1} \OperatorTok{+} \VariableTok{X2}\OperatorTok{;}   \CommentTok{\% Combined noise}

\VariableTok{fprintf}\NormalTok{(}\SpecialStringTok{\textquotesingle{}Var(X1) = \%.4f (theory: \%.4f)\textbackslash{}n\textquotesingle{}}\OperatorTok{,} \VariableTok{var}\NormalTok{(}\VariableTok{X1}\NormalTok{)}\OperatorTok{,} \VariableTok{sigma1}\OperatorTok{\^{}}\FloatTok{2}\NormalTok{)}\OperatorTok{;}
\VariableTok{fprintf}\NormalTok{(}\SpecialStringTok{\textquotesingle{}Var(X2) = \%.4f (theory: \%.4f)\textbackslash{}n\textquotesingle{}}\OperatorTok{,} \VariableTok{var}\NormalTok{(}\VariableTok{X2}\NormalTok{)}\OperatorTok{,} \VariableTok{sigma2}\OperatorTok{\^{}}\FloatTok{2}\NormalTok{)}\OperatorTok{;}
\VariableTok{fprintf}\NormalTok{(}\SpecialStringTok{\textquotesingle{}Var(X1+X2) = \%.4f (theory: \%.4f)\textbackslash{}n\textquotesingle{}}\OperatorTok{,} \VariableTok{var}\NormalTok{(}\VariableTok{Y}\NormalTok{)}\OperatorTok{,} \VariableTok{sigma1}\OperatorTok{\^{}}\FloatTok{2} \OperatorTok{+} \VariableTok{sigma2}\OperatorTok{\^{}}\FloatTok{2}\NormalTok{)}\OperatorTok{;}
\VariableTok{fprintf}\NormalTok{(}\SpecialStringTok{\textquotesingle{}Std(X1+X2) = \%.4f (theory: \%.4f)\textbackslash{}n\textquotesingle{}}\OperatorTok{,} \VariableTok{std}\NormalTok{(}\VariableTok{Y}\NormalTok{)}\OperatorTok{,} \VariableTok{sqrt}\NormalTok{(}\VariableTok{sigma1}\OperatorTok{\^{}}\FloatTok{2}\OperatorTok{+}\VariableTok{sigma2}\OperatorTok{\^{}}\FloatTok{2}\NormalTok{))}\OperatorTok{;}
\VariableTok{fprintf}\NormalTok{(}\SpecialStringTok{\textquotesingle{}\textbackslash{}nNote: std devs do NOT add: \%.3f + \%.3f = \%.3f ≠ \%.3f\textbackslash{}n\textquotesingle{}}\OperatorTok{,} \OperatorTok{...}
    \VariableTok{sigma1}\OperatorTok{,} \VariableTok{sigma2}\OperatorTok{,} \VariableTok{sigma1}\OperatorTok{+}\VariableTok{sigma2}\OperatorTok{,} \VariableTok{sqrt}\NormalTok{(}\VariableTok{sigma1}\OperatorTok{\^{}}\FloatTok{2}\OperatorTok{+}\VariableTok{sigma2}\OperatorTok{\^{}}\FloatTok{2}\NormalTok{))}\OperatorTok{;}
\end{Highlighting}
\end{Shaded}

\end{PSbox}

\begin{PSbox}[unbreakable]{ex}{Example 3: Mean and Variance of Transformed RV}

\tcbset{PScodemode/.style={unbreakable}}

\begin{Shaded}
\begin{Highlighting}[]
\CommentTok{\%\% E[g(X)] using LOTUS: Power from Voltage}
\CommentTok{\% X = voltage \textasciitilde{} N(0, 0.5V)}
\CommentTok{\% P = X\^{}2 / R, with R = 50 ohms}

\VariableTok{sigma} \OperatorTok{=} \FloatTok{0.5}\OperatorTok{;}    \CommentTok{\% noise voltage RMS}
\VariableTok{R} \OperatorTok{=} \FloatTok{50}\OperatorTok{;}         \CommentTok{\% resistance}
\VariableTok{N} \OperatorTok{=} \FloatTok{100000}\OperatorTok{;}

\VariableTok{X} \OperatorTok{=} \VariableTok{sigma} \OperatorTok{*} \VariableTok{randn}\NormalTok{(}\FloatTok{1}\OperatorTok{,} \VariableTok{N}\NormalTok{)}\OperatorTok{;}
\VariableTok{P} \OperatorTok{=} \VariableTok{X}\OperatorTok{.\^{}}\FloatTok{2} \OperatorTok{/} \VariableTok{R}\OperatorTok{;}   \CommentTok{\% instantaneous power}

\CommentTok{\% Using LOTUS: E[X\^{}2/R] = E[X\^{}2]/R = sigma\^{}2/R}
\VariableTok{E\_P\_theory} \OperatorTok{=} \VariableTok{sigma}\OperatorTok{\^{}}\FloatTok{2} \OperatorTok{/} \VariableTok{R}\OperatorTok{;}
\VariableTok{E\_P\_sim} \OperatorTok{=} \VariableTok{mean}\NormalTok{(}\VariableTok{P}\NormalTok{)}\OperatorTok{;}

\VariableTok{fprintf}\NormalTok{(}\SpecialStringTok{\textquotesingle{}Mean power: Theory = \%.6f W, Sim = \%.6f W\textbackslash{}n\textquotesingle{}}\OperatorTok{,} \VariableTok{E\_P\_theory}\OperatorTok{,} \VariableTok{E\_P\_sim}\NormalTok{)}\OperatorTok{;}
\VariableTok{fprintf}\NormalTok{(}\SpecialStringTok{\textquotesingle{}Mean power = \%.3f mW\textbackslash{}n\textquotesingle{}}\OperatorTok{,} \VariableTok{E\_P\_theory} \OperatorTok{*} \FloatTok{1000}\NormalTok{)}\OperatorTok{;}
\end{Highlighting}
\end{Shaded}

\end{PSbox}

\begin{PSbox}[unbreakable]{ex}{Example 4: Higher Moments and Skewness}

\tcbset{PScodemode/.style={unbreakable}}

\begin{Shaded}
\begin{Highlighting}[]
\CommentTok{\%\% Skewness and Kurtosis Comparison}
\VariableTok{N} \OperatorTok{=} \FloatTok{100000}\OperatorTok{;}

\CommentTok{\% Gaussian (symmetric)}
\VariableTok{X\_gauss} \OperatorTok{=} \VariableTok{randn}\NormalTok{(}\FloatTok{1}\OperatorTok{,} \VariableTok{N}\NormalTok{)}\OperatorTok{;}

\CommentTok{\% Exponential (right{-}skewed)}
\VariableTok{X\_exp} \OperatorTok{=} \VariableTok{exprnd}\NormalTok{(}\FloatTok{1}\OperatorTok{,} \FloatTok{1}\OperatorTok{,} \VariableTok{N}\NormalTok{)}\OperatorTok{;}

\CommentTok{\% Uniform (symmetric, light tails)}
\VariableTok{X\_unif} \OperatorTok{=} \VariableTok{rand}\NormalTok{(}\FloatTok{1}\OperatorTok{,} \VariableTok{N}\NormalTok{)}\OperatorTok{;}

\VariableTok{fprintf}\NormalTok{(}\SpecialStringTok{\textquotesingle{}Distribution    | Skewness | Kurtosis\textbackslash{}n\textquotesingle{}}\NormalTok{)}\OperatorTok{;}
\VariableTok{fprintf}\NormalTok{(}\SpecialStringTok{\textquotesingle{}{-}{-}{-}{-}{-}{-}{-}{-}{-}{-}{-}{-}{-}{-}{-}{-}+{-}{-}{-}{-}{-}{-}{-}{-}{-}{-}+{-}{-}{-}{-}{-}{-}{-}{-}{-}\textbackslash{}n\textquotesingle{}}\NormalTok{)}\OperatorTok{;}
\VariableTok{fprintf}\NormalTok{(}\SpecialStringTok{\textquotesingle{}Gaussian        | \%+.3f   | \%+.3f\textbackslash{}n\textquotesingle{}}\OperatorTok{,} \VariableTok{skewness}\NormalTok{(}\VariableTok{X\_gauss}\NormalTok{)}\OperatorTok{,} \VariableTok{kurtosis}\NormalTok{(}\VariableTok{X\_gauss}\NormalTok{)}\OperatorTok{{-}}\FloatTok{3}\NormalTok{)}\OperatorTok{;}
\VariableTok{fprintf}\NormalTok{(}\SpecialStringTok{\textquotesingle{}Exponential     | \%+.3f   | \%+.3f\textbackslash{}n\textquotesingle{}}\OperatorTok{,} \VariableTok{skewness}\NormalTok{(}\VariableTok{X\_exp}\NormalTok{)}\OperatorTok{,} \VariableTok{kurtosis}\NormalTok{(}\VariableTok{X\_exp}\NormalTok{)}\OperatorTok{{-}}\FloatTok{3}\NormalTok{)}\OperatorTok{;}
\VariableTok{fprintf}\NormalTok{(}\SpecialStringTok{\textquotesingle{}Uniform         | \%+.3f   | \%+.3f\textbackslash{}n\textquotesingle{}}\OperatorTok{,} \VariableTok{skewness}\NormalTok{(}\VariableTok{X\_unif}\NormalTok{)}\OperatorTok{,} \VariableTok{kurtosis}\NormalTok{(}\VariableTok{X\_unif}\NormalTok{)}\OperatorTok{{-}}\FloatTok{3}\NormalTok{)}\OperatorTok{;}
\VariableTok{fprintf}\NormalTok{(}\SpecialStringTok{\textquotesingle{}\textbackslash{}nTheoretical:\textbackslash{}n\textquotesingle{}}\NormalTok{)}\OperatorTok{;}
\VariableTok{fprintf}\NormalTok{(}\SpecialStringTok{\textquotesingle{}Gaussian:    skew=0, kurt=0\textbackslash{}n\textquotesingle{}}\NormalTok{)}\OperatorTok{;}
\VariableTok{fprintf}\NormalTok{(}\SpecialStringTok{\textquotesingle{}Exponential: skew=2, kurt=6\textbackslash{}n\textquotesingle{}}\NormalTok{)}\OperatorTok{;}
\VariableTok{fprintf}\NormalTok{(}\SpecialStringTok{\textquotesingle{}Uniform:     skew=0, kurt={-}1.2\textbackslash{}n\textquotesingle{}}\NormalTok{)}\OperatorTok{;}
\end{Highlighting}
\end{Shaded}

\end{PSbox}

\begin{PSbox}[unbreakable]{ex}{Example 5: MGF Verification}

\tcbset{PScodemode/.style={unbreakable}}

\begin{Shaded}
\begin{Highlighting}[]
\CommentTok{\%\% Verify MGF: Compute moments by differentiation (numerical)}
\CommentTok{\% For X \textasciitilde{} Exponential(lambda)}
\VariableTok{lambda\_rate} \OperatorTok{=} \FloatTok{2}\OperatorTok{;}
\VariableTok{beta\_scale} \OperatorTok{=} \FloatTok{1}\OperatorTok{/}\VariableTok{lambda\_rate}\OperatorTok{;}

\CommentTok{\% MGF: M(t) = lambda/(lambda {-} t)}
\CommentTok{\% M\textquotesingle{}(0) = E[X], M\textquotesingle{}\textquotesingle{}(0) = E[X\^{}2]}

\VariableTok{t} \OperatorTok{=} \FloatTok{0}\OperatorTok{;}
\VariableTok{dt} \OperatorTok{=} \FloatTok{1e{-}6}\OperatorTok{;}

\CommentTok{\% Numerical derivatives}
\VariableTok{M} \OperatorTok{=} \OperatorTok{@}\NormalTok{(}\VariableTok{t}\NormalTok{) }\VariableTok{lambda\_rate} \OperatorTok{./}\NormalTok{ (}\VariableTok{lambda\_rate} \OperatorTok{{-}} \VariableTok{t}\NormalTok{)}\OperatorTok{;}
\VariableTok{M\_prime} \OperatorTok{=}\NormalTok{ (}\VariableTok{M}\NormalTok{(}\VariableTok{t}\OperatorTok{+}\VariableTok{dt}\NormalTok{) }\OperatorTok{{-}} \VariableTok{M}\NormalTok{(}\VariableTok{t}\OperatorTok{{-}}\VariableTok{dt}\NormalTok{)) }\OperatorTok{/}\NormalTok{ (}\FloatTok{2}\OperatorTok{*}\VariableTok{dt}\NormalTok{)}\OperatorTok{;}
\VariableTok{M\_double\_prime} \OperatorTok{=}\NormalTok{ (}\VariableTok{M}\NormalTok{(}\VariableTok{t}\OperatorTok{+}\VariableTok{dt}\NormalTok{) }\OperatorTok{{-}} \FloatTok{2}\OperatorTok{*}\VariableTok{M}\NormalTok{(}\VariableTok{t}\NormalTok{) }\OperatorTok{+} \VariableTok{M}\NormalTok{(}\VariableTok{t}\OperatorTok{{-}}\VariableTok{dt}\NormalTok{)) }\OperatorTok{/} \VariableTok{dt}\OperatorTok{\^{}}\FloatTok{2}\OperatorTok{;}

\VariableTok{fprintf}\NormalTok{(}\SpecialStringTok{\textquotesingle{}From MGF derivatives:\textbackslash{}n\textquotesingle{}}\NormalTok{)}\OperatorTok{;}
\VariableTok{fprintf}\NormalTok{(}\SpecialStringTok{\textquotesingle{}  E[X] = M\textquotesingle{}\textquotesingle{}(0) = \%.4f (theory: \%.4f)\textbackslash{}n\textquotesingle{}}\OperatorTok{,} \VariableTok{M\_prime}\OperatorTok{,} \FloatTok{1}\OperatorTok{/}\VariableTok{lambda\_rate}\NormalTok{)}\OperatorTok{;}
\VariableTok{fprintf}\NormalTok{(}\SpecialStringTok{\textquotesingle{}  E[X²] = M\textquotesingle{}\textquotesingle{}\textquotesingle{}\textquotesingle{}(0) = \%.4f (theory: \%.4f)\textbackslash{}n\textquotesingle{}}\OperatorTok{,} \VariableTok{M\_double\_prime}\OperatorTok{,} \FloatTok{2}\OperatorTok{/}\VariableTok{lambda\_rate}\OperatorTok{\^{}}\FloatTok{2}\NormalTok{)}\OperatorTok{;}
\VariableTok{fprintf}\NormalTok{(}\SpecialStringTok{\textquotesingle{}  Var(X) = E[X²] {-} (E[X])² = \%.4f (theory: \%.4f)\textbackslash{}n\textquotesingle{}}\OperatorTok{,} \OperatorTok{...}
    \VariableTok{M\_double\_prime} \OperatorTok{{-}} \VariableTok{M\_prime}\OperatorTok{\^{}}\FloatTok{2}\OperatorTok{,} \FloatTok{1}\OperatorTok{/}\VariableTok{lambda\_rate}\OperatorTok{\^{}}\FloatTok{2}\NormalTok{)}\OperatorTok{;}
\end{Highlighting}
\end{Shaded}

\end{PSbox}

\PSsection{4.12}{Practice Problems}

\begin{PSbox}[unbreakable]{prob}{Problem 1: Basic Computation}

A discrete RV $X$ has PMF: $P(X = -1) = 0.2,\, P(X = 0) = 0.5,\, P(X = 1) = 0.2,\, P(X = 2) = 0.1$.

\begin{enumerate}
\def\labelenumi{(\alph{enumi})}
\tightlist
\item
  Find $E[X]$.
\item
  Find $E[X^{2}]$.
\item
  Find $\operatorname{Var}(X)$.
\item
  Find $E[3X + 2]$.
\item
  Find $\operatorname{Var}(3X + 2)$.
\end{enumerate}

\PSsolution{} 

\begin{enumerate}
\def\labelenumi{(\alph{enumi})}
\item
  $E[X] = (-1)(0.2) + (0)(0.5) + (1)(0.2) + (2)(0.1) = -0.2 + 0 + 0.2 + 0.2 = 0.2$
\item
  $E[X^{2}] = (1)(0.2) + (0)(0.5) + (1)(0.2) + (4)(0.1) = 0.2 + 0 + 0.2 + 0.4 = 0.8$
\item
  $\operatorname{Var}(X) = E[X^{2}] - (E[X])^{2} = 0.8 - 0.04 = 0.76$
\item
  $E[3X + 2] = 3E[X] + 2 = 3(0.2) + 2 = 2.6$
\item
  $\operatorname{Var}(3X + 2) = 9 \cdot \operatorname{Var}(X) = 9(0.76) = 6.84$
\end{enumerate}

\end{PSbox}

\begin{PSbox}[unbreakable]{prob}{Problem 2: Signal and Noise}

A received signal $R = 2.5 + N$, where $N \sim N(0,\, 0.04)$.

\begin{enumerate}
\def\labelenumi{(\alph{enumi})}
\tightlist
\item
  Find $E[R]$ and $\operatorname{Var}(R)$.
\item
  Find the SNR in dB.
\item
  Find $P(R < 0)$ — probability the signal goes negative.
\end{enumerate}

\PSsolution{} 

\begin{enumerate}
\def\labelenumi{(\alph{enumi})}
\item
  $E[R] = 2.5,\, \operatorname{Var}(R) = \operatorname{Var}(N) = 0.04,\, \sigma_{R} = 0.2$
\item
  Signal power $= 2.5^{2} = 6.25$, Noise power = 0.04\PSbr{}$\mathrm{SNR} = \frac{6.25}{0.04} = 156.25 \to \mathrm{SNR}_{dB} = 10 \cdot \log_{10}(156.25) = 21.94$ dB
\item
  $P(R < 0) = P(N < -2.5) = \Phi \left(-\frac{2.5}{0.2}\right) = \Phi (-12.5) \approx 0$ (extremely unlikely)
\end{enumerate}

\end{PSbox}

\begin{PSbox}[unbreakable]{prob}{Problem 3: Averaging Measurements}

A sensor makes $n$ independent measurements of a constant voltage $V = 3.3\,\mathrm{V}$.\PSbr{}Each measurement: $X_{i} = 3.3 + N_{i}$ where $N_{i} \sim N(0,\, 0.01)$.

\begin{enumerate}
\def\labelenumi{(\alph{enumi})}
\tightlist
\item
  What is $E[\bar{X}]$ and $\operatorname{Var}(\bar{X})$ for $n = 1$?
\item
  How many measurements $n$ needed so that $\sigma_{\bar{X}} < 0.01\,\mathrm{V}$?
\item
  With $n = 100$, what is $P(\lvert \bar{X} - 3.3\rvert > 0.02)$?
\end{enumerate}

\PSsolution{} 

\begin{enumerate}
\def\labelenumi{(\alph{enumi})}
\item
  $\displaystyle E[\bar{X}] = 3.3,\, \operatorname{Var}(\bar{X}) = \frac{0.01}{1} = 0.01,\, \sigma = 0.1\,\mathrm{V}$
\item
  $\sigma_{\bar{X}} = \frac{\sigma}{\sqrt{n}} < 0.01 \to \sqrt{n} > 10 \to n > 100$ \ensuremath{\to} need $n \ge 101$
\item
  $\displaystyle \sigma_{\bar{X}} = \frac{0.1}{\sqrt{100}} = 0.01\,\mathrm{V}$\PSbr{}$P(\lvert \bar{X} - 3.3\rvert > 0.02) = P(\lvert Z\rvert > 2) = 2(1-\Phi (2)) = 0.0455$
\end{enumerate}

\end{PSbox}

\begin{PSbox}[unbreakable]{prob}{Problem 4: Component Lifetime}

Components have lifetime $T \sim \mathrm{Exp}(\lambda = 0.002\ \text{per hour})$.

\begin{enumerate}
\def\labelenumi{(\alph{enumi})}
\tightlist
\item
  Find MTTF, $\sigma_{T}$, and coefficient of variation.
\item
  Find $P(T > 1000\,\mathrm{hours})$.
\item
  Find $E[T^{2}]$ and interpret.
\end{enumerate}

\PSsolution{} 

\begin{enumerate}
\def\labelenumi{(\alph{enumi})}
\item
  $\mathrm{MTTF} = \frac{1}{\lambda} = 500$ hours, $\sigma_{T} = 500$ hours, CV $= \frac{\sigma}{\mu} = 1$
\item
  $P(T > 1000) = e^{-0.002 \times 1000} = e^{-2} = 0.1353$
\item
  $E[T^{2}] = \operatorname{Var}(T) + (E[T])^{2} = 500^{2} + 500^{2} = 500000\,\mathrm{hours}^{2}$\PSbr{}$\sqrt{E[T^{2}]} = 707$ hours (RMS lifetime)
\end{enumerate}

\end{PSbox}

\begin{PSbox}[unbreakable]{prob}{Problem 5: MGF Application}

$X$ has MGF $M_{X}(t) = \frac{1}{3}e^{t} + \frac{2}{3}e^{2t}$.

\begin{enumerate}
\def\labelenumi{(\alph{enumi})}
\tightlist
\item
  What type of RV is $X$? Find its PMF.
\item
  Find $E[X]$ using the MGF.
\item
  Find $\operatorname{Var}(X)$ using the MGF.
\end{enumerate}

\PSsolution{} 

\begin{enumerate}
\def\labelenumi{(\alph{enumi})}
\item
  $X$ is discrete. Comparing with $\sum p_{X}(x)e^{tx}$:\PSbr{}$\displaystyle P(X = 1) = \frac{1}{3},\, P(X = 2) = \frac{2}{3}$
\item
  $\displaystyle M'(t) = \frac{1}{3}e^{t} + \frac{4}{3}e^{2t}$\PSbr{}$\displaystyle E[X] = M'(0) = \frac{1}{3} + \frac{4}{3} = \frac{5}{3}$
\item
  $\displaystyle M''(t) = \frac{1}{3}e^{t} + \frac{8}{3}e^{2t}$\PSbr{}$\displaystyle E[X^{2}] = M''(0) = \frac{1}{3} + \frac{8}{3} = 3$\PSbr{}$\displaystyle \operatorname{Var}(X) = 3 - \frac{5}{3}^{2} = 3 - \frac{25}{9} = \frac{2}{9}$
\end{enumerate}

\end{PSbox}

\PSsection{4.13}{Key Takeaways}

\begin{enumerate}
\def\labelenumi{\arabic{enumi}.}
\tightlist
\item
  \textbf{$E[X]$ = long-run average} — the center of the distribution, DC component
\item
  \textbf{$\operatorname{Var}(X)$ = spread/power} — for zero-mean signals, variance IS power
\item
  \textbf{$\sigma$ has same units as $X$} — directly interpretable as RMS value
\item
  \textbf{Linearity of expectation} always holds; variance addition requires independence
\item
  \textbf{MGF uniquely identifies} a distribution and generates all moments
\item
  \textbf{In engineering:} mean \ensuremath{\to} signal level, variance \ensuremath{\to} noise power, std dev \ensuremath{\to} RMS noise
\end{enumerate}

\emph{Next Module: {Module 5 — Important Random-Variable Functions}}
